\section{Data Source and Collection}
\label{sec:3-1}

The project uses the publicly available \textbf{Online Retail II} dataset as a historical extract of e-commerce transaction records. The dataset is obtained and loaded into the analytical environment for data inspection, preprocessing and customer churn analysis. It contains transaction-level information, including invoice identifiers, customer identifiers, product codes, quantities, prices and invoice dates.

\section{Dataset Characteristics and Limitations}
\label{sec:3-2}

The dataset covers the period from December 2009 to December 2011. Its relatively limited historical coverage constrains the number of customer snapshots and future observation periods available for churn labeling (Section~\ref{sec:5-9}). In addition, the retailer serves both individual customers and wholesalers, resulting in substantial differences in purchasing frequency and transaction value. Seasonal demand, particularly around Christmas, may also influence purchasing patterns and should be considered when designing temporal validation and evaluating model performance.

\section{Raw Data Schema}
\label{sec:3-3}

The raw dataset is a single flat transaction table (not pre-split into separate customer, order and product tables): each row is one product line within one invoice.

\begin{center}
\begin{tabular}{@{}p{2.6cm}p{2cm}p{10cm}@{}}
    \toprule
    Column & Type & Description \\
    \midrule
    \texttt{Invoice} & string & Invoice number. A ``C''-prefixed value (e.g.\ \texttt{C536379}) marks a cancellation of an earlier invoice; \texttt{Quantity} is negative on these rows. \\
    \texttt{StockCode} & string & Product code. Most codes are numeric (e.g.\ \texttt{85123A}); a small set of non-numeric codes are not products at all (below). \\
    \texttt{Description} & string & Free-text product description; can be missing or inconsistent for the same \texttt{StockCode}. \\
    \texttt{Quantity} & integer & Units on this line. Negative for cancellations. \\
    \texttt{InvoiceDate} & datetime & Timestamp of the transaction line. \\
    \texttt{Price} & float & Unit price (GBP). \\
    \texttt{Customer ID} & float & Customer identifier. Missing for a meaningful share of rows (Section~\ref{sec:5-1}) -- these cannot be attached to any customer history. \\
    \texttt{Country} & string & Customer's country, self-reported at checkout. \\
    \bottomrule
\end{tabular}
\end{center}

\section{Pipeline Column Names}
\label{sec:3-4}

These raw columns are renamed to the snake\_case names used throughout the pipeline (code, notebooks and the rest of this report), and two columns are derived from them:

\begin{center}
\begin{tabular}{@{}p{3cm}p{4cm}p{7.6cm}@{}}
    \toprule
    Raw column & Pipeline name & Notes \\
    \midrule
    \texttt{Invoice} & \texttt{invoice\_id} & -- \\
    \texttt{StockCode} & \texttt{product\_id} & -- \\
    \texttt{Description} & \texttt{product\_description} & -- \\
    \texttt{Quantity} & \texttt{quantity} & -- \\
    \texttt{InvoiceDate} & \texttt{invoice\_time} & -- \\
    \texttt{Price} & \texttt{unit\_price} & -- \\
    \texttt{Customer ID} & \texttt{customer\_id} & -- \\
    \texttt{Country} & \texttt{country} & -- \\
    \emph{(derived)} & \texttt{total\_price} & $= \texttt{quantity} \times \texttt{unit\_price}$ \\
    \emph{(derived)} & \texttt{invoice\_month} & \texttt{invoice\_time} truncated to the first of its month, used for monthly aggregation \\
    \bottomrule
\end{tabular}
\end{center}

\section{Value Conventions}
\label{sec:3-5}

Three value conventions recur throughout this report and are worth stating once, up front:

\begin{itemize}
    \item \textbf{Cancellation invoices.} An \texttt{invoice\_id} starting with ``C'' marks a cancellation of a previous purchase, with negative \texttt{quantity}. These rows are excluded from purchase history (Section~\ref{sec:5-2}) but kept separately to build the cancellation features in Section~\ref{sec:5-6}.
    \item \textbf{Non-product stock codes.} A small number of \texttt{product\_id} values are not real products but operational/accounting entries -- postage (\texttt{POST}, \texttt{DOT}), manual adjustments (\texttt{M}), carriage (\texttt{C2}), discounts (\texttt{D}), samples (\texttt{S}), bank charges and similar. These are discovered empirically (Section~\ref{sec:5-1}), not assumed in advance, and are excluded from product- and customer-level analysis.
    \item \textbf{Missing Customer ID.} Rows without a \texttt{customer\_id} cannot be attached to any customer history and are excluded from every customer-level analysis from Section~\ref{sec:4-3} onward, although they are still counted in the raw totals of Section~\ref{sec:4-2}.
\end{itemize}

\section{Information Not Available}
\label{sec:3-6}

The dataset contains transactions only. Product cost, product category hierarchy and margin are not available, and neither is any marketing or campaign information (campaign dates, channels, treatment/control assignment, discounts, voucher or contact costs, campaign outcomes). A production system would ideally collect these; in this report, the related quantities used for business-impact estimation in Chapter~\ref{ch:evaluation} (margin, contact cost, incentive cost, retention effectiveness) are explicit assumptions.

\section{Analytical Dataset Design}
\label{sec:3-7}

Transactions are not modeled directly. Each final modeling row represents a \textbf{Customer $\times$ Snapshot Date}: the customer's historical features computed only from information available up to the snapshot date, together with a future churn label. This dataset is built in Chapter~\ref{ch:data-preparation}.
